\documentclass{article}
\usepackage[utf8]{inputenc}
\usepackage[T1]{fontenc}
\usepackage[swedish]{babel}
\usepackage{graphicx}
\usepackage{enumitem}
\usepackage{fancyhdr}
\usepackage[table]{xcolor}
\usepackage{tabularx}
\usepackage{changepage}
\usepackage[most]{tcolorbox}
\usepackage{amsmath}
\usepackage{csquotes}
\usepackage{hyperref}
\usepackage{url}

\hypersetup{
    colorlinks=true,
    linkcolor=blue,
    urlcolor=blue,
    pdftitle={Metodbeskrivning för riskbedömning av kund}
}
\title{Metodbeskrivning för riskbedömning av kund\\\large Version 1.0}
\pagestyle{fancy}
\fancyhf{}
\fancyhead[R]{Metodbeskrivning för riskbedömning av kund, \today}

\begin{document}



\vspace{1em}

\maketitle
\tableofcontents
\newpage

\section*{Inledning}
\addcontentsline{toc}{section}{Inledning}

Denna metodbeskrivning fastställer hur vi som redovisningsbyrå genomför riskbedömningen av kunder, både vid antagandet av nya uppdrag och vid den löpande uppdateringen av befintliga kunders riskprofiler. Syftet är att säkerställa en enhetlig och rättssäker tillämpning av penningtvättslagstiftningens krav på kundkännedom och riskklassificering.

För att upprätthålla ett aktuellt och rättssäkert arbetssätt uppdateras detta dokument minst en gång per år, eller vid behov till följd av ändrad lagstiftning, vägledning från tillsynsmyndigheter eller interna riskobservationer. Ansvarig för uppdatering och godkännande är den verksamhetsansvarige. Vid varje ändring uppdateras dokumentets versionsnummer.

% Definitioner: hot och sårbarhet
\section*{Definition av risknivåer}
\addcontentsline{toc}{section}{Definition av risknivåer}

För att kunna uppfylla penningtvättslagens krav på riskbaserade åtgärder (2 kap. 1 § och 3 §) krävs en systematisk och dokumenterad metod för att identifiera och bedöma risker. 

Vi inleder därför med en definition av de risknivåer som tillämpas i vår verksamhet. Risknivån utgör en kombination av externa hot (kundrisk) och interna sårbarheter (tjänstens känslighet), vilket ligger till grund för beslut om kontrollåtgärder och kundacceptans.

\begin{itemize}[label=--]
  \item \textbf{Hot} – den yttre risken: i vilken grad en viss kundtyp eller verksamhet kan utnyttjas för penningtvätt eller finansiering av terrorism. Bedömningen utgår från den nationella riskbedömningen och avser hur attraktiv kunden är för kriminella aktörer.

  \item \textbf{Sårbarhet} – den inre risken: i vilken utsträckning våra egna tjänster kan manipuleras eller missbrukas utan att det upptäcks. Faktorer inkluderar dokumentation, kontrollrutiner, spårbarhet och insyn.
\end{itemize}

\subsection*{Princip för riskbedömning}
\addcontentsline{toc}{subsection}{Princip för riskbedömning}

Vi tillämpar en enkel och transparent modell där:

\begin{quote}
\textbf{Risknivå = Hot (kund) × Sårbarhet (tjänst)}
\end{quote}

Hotnivån hämtas från den \textit{Nationella riskbedömningen av penningtvätt och finansiering av terrorism i Sverige 2020/2021}\footnote{\url{https://www.ekobrottsmyndigheten.se/wp-content/uploads/2021/05/nationell-riskbedomning-av-penningtvatt-och-finansiering-av-terrorism-2020-2021.pdf}}. Sårbarhetsnivån fastställs internt för varje tjänst utifrån dess kontrollbarhet.

\textbf{Syftet är att styra vilka åtgärder som krävs enligt penningtvättslagen, med riskpoängen som beslutsunderlag.}

%%%%%%%%%%%%%%%%%%%%%%%%%%%

\vspace{3em}

%% Allmänt Hot x Sårbarhet

\section*{Metod för allmän riskbedömning}
\addcontentsline{toc}{section}{Metod för allmän riskbedömning}

Vår metod för allmän riskbedömning baseras på den riskmatris och skala som presenteras i \textit{Nationell riskbedömning av penningtvätt och finansiering av terrorism 2020/2021}, utgiven av bland andra Finansdepartementet, Polismyndigheten, Ekobrottsmyndigheten och Finansinspektionen.

Risknivån beräknas genom att multiplicera en bedömd hotnivå med en bedömd sårbarhetsnivå enligt följande formel:

\[
\text{Risknivå} = \text{Hotnivå} \times \text{Sårbarhetsnivå}
\]

Varje komponent bedöms på en fyrgradig skala:

\begin{itemize}[label=--]
  \item 1 = Låg
  \item 2 = Medel
  \item 3 = Betydande
  \item 4 = Hög
\end{itemize}

\noindent Skalan följer den semikvantitativa metodik som används i den nationella riskbedömningen (bilaga ”Process och metod”, s. 123 ff.). Genom att använda denna metod säkerställs att vår klassificering är spårbar, jämförbar och förankrad i myndighetsgemensamma riktlinjer.

För varje kombination av kundtyp och tjänst motiverar vi bedömda siffror i separata tabeller i efterföljande avsnitt.

%%%%%%%%%%%%%%%%%%%%%%%%%%%%%
\vspace{3em}

%% Hot
%\section*{Motivering av hotnivåer per kundtyp}
%\addcontentsline{toc}{section}{Motivering av hotnivåer per kundtyp}


\subsection*{Härledning av hotnivåer per verksamhetsområde enligt myndighetskatalog}
\addcontentsline{toc}{subsection}{Härledning av hotnivåer per verksamhetsområde enligt myndighetskatalog}


Som underlag för vår generella riskbedömning har vi utgått från den myndighetsgemensamma rapporten \textit{Nationell riskbedömning av penningtvätt och finansiering av terrorism i Sverige 2020/2021}, utgiven av Ekobrottsmyndigheten, Polismyndigheten, Finansinspektionen m.fl.

Dokumentet innehåller i kapitel 7 (s. 39–117) en sektorskatalog där samtliga verksamhetsområden som omfattas av penningtvättslagstiftningen bedöms individuellt utifrån risken att utnyttjas för penningtvätt eller finansiering av terrorism. För varje sektor anges en risknivå i termerna \textit{låg}, \textit{medel}, \textit{betydande} eller \textit{hög}, med tillhörande motivering.

Följande tabell återger dessa myndighetsbedömningar, som vi använder som referens vid klassificering av våra kunders branschtillhörighet:

\begin{center}
\begin{tabular}{|p{9cm}|c|}
\hline
\textbf{Verksamhetsområde} & \textbf{Hotnivå} \\
\hline
Bank- eller finansieringsrörelse & Betydande \\
Livförsäkringsrörelse & Medel \\
Värdepappersrörelse & Medel \\
Valutaväxling, finansiell verksamhet & Hög \\
Försäkringsförmedlare & Medel\\
Utgivare av elektroniska pengar & Betydande \\
Fondförvaltning & Medel\\
Betalningsinstitut & Betydande \\
Registrerade betaltjänstleverantörer & Betydande \\
Verksamhet med konsumentkrediter & Medel \\
Verksamhet med bostadskrediter & Medel \\
Fastighetsmäklare med fullständig registrering & Betydande \\
Varuhandlare (inkl. e-handel i tillämpliga fall) & Hög \\
Pantbank & Medel \\
Bokförings- och revisionstjänster & Betydande \\
Skatterådgivare & Betydande \\
Oberoende jurist & Betydande \\
Bolagsbildare och företagsmäklare & Hög \\
Kontorshotell och postboxföretag & Betydande \\
Trustförvaltare & Hög \\
Styrelserepresentation, nominella aktieägare & Medel \\
Verksamhet som auktoriserad och godkänd revisor & Medel \\
Advokater och advokatbyråer & Medel \\
Svenska spelmarknaden & Hög \\
\hline
\end{tabular}
\end{center}

De ovanstående nivåerna återspeglar den kvalitativa skala som används i den nationella riskbedömningen. Noterbart är att \textit{betydande} utgör en separat riskkategori mellan \textit{medel} och \textit{hög} och bör hanteras som ett distinkt värde vid kvantifiering.

Den fullständiga rapporten finns tillgänglig via Ekobrottsmyndighetens webbplats och vi hänvisar uttryckligen till denna källa för att säkerställa spårbarhet, förankring och legitimitet i vår klassificeringsmodell.

\subsubsection*{Översättning av hotnivåer till numerisk skala}

För att möjliggöra en enhetlig och kvantitativ tillämpning av riskbedömningen har vi översatt myndigheternas verbala riskkategorier till en numerisk skala enligt följande:

\begin{center}
\begin{tabular}{|c|c|}
\hline
\textbf{Benämning i myndighetsrapport} & \textbf{Numeriskt värde i vår modell} \\
\hline
Låg & 1 \\
Medel & 2 \\
Betydande & 3 \\
Hög & 4 \\
\hline
\end{tabular}
\end{center}

Vi tillämpar denna fyrgradiga skala både för hotnivå (extern exponering) och sårbarhetsnivå (intern kontrollförmåga). Den resulterande risknivån erhålls genom multiplikation:

\[
\text{Riskvärde} = \text{Hotnivå} \times \text{Sårbarhetsnivå}
\]

Detta ger ett slutligt intervall från 1 till 16 och ligger till grund för våra prioriteringar och åtgärdsplaner.

\begin{tcolorbox}
\textbf{Exempel:} En restaurangverksamhet med hotnivå 4 och en tillhandahållen tjänst med sårbarhetsnivå 2 får ett sammanlagt riskvärde på \(4 \times 2 = 8\).
\end{tcolorbox}


% Exempel på bedömning av hotnivåer

\subsubsection*{Kommunalt fastighetsbolag – Hotnivå 1 (låg)}

Kunden är en kommunalt ägd juridisk person och betraktas därmed som offentlig aktör. Sådana aktörer omfattas av offentlighetsprincipen, interna revisionskrav samt upphandlingslagstiftning, vilket tillsammans medför ett högt mått av transparens och dokumentationsplikt. 

Enligt den nationella riskbedömningen (s. 17 och 38) bedöms risken för att offentliga verksamheter används för penningtvätt generellt som låg. Det saknas i regel kommersiella incitament för att möjliggöra eller dölja olagliga transaktioner, och kontrollsystemen är robusta. 

Mot denna bakgrund tilldelas hotnivån för denna kundkategori värdet \textbf{1 – låg extern risk}.




\subsubsection*{Bygg-, städ-, restaurang- och transportbranschen – Hotnivå 4 (hög)}

Den \textit{Nationella riskbedömningen av penningtvätt och finansiering av terrorism i Sverige 2020/2021} konstaterar följande (s. 15):

\begin{quote}
\itshape
Stora belopp kommer traditionellt sett bland annat från skattebrott. Dessa är till stor del relaterade till oredovisade skatter och avgifter gällande arbetskraft inom främst bygg-, städ- och transportsektorn.
\end{quote}

Vidare anges (kap. 7.15, s. 91):

\begin{quote}
\itshape
Det är vanligt att [redovisningsföretag] har kunder inom branscher som är kända för att vara mer utsatta för ekonomisk brottslighet än andra, till exempel inom bygg- och restaurangbranschen. [...] Möjligheten att bedriva penningtvätt i sektorn bedöms som hög.
\end{quote}

Utifrån detta bedömer vi följande branscher som särskilt exponerade för hot kopplade till penningtvätt:

\begin{itemize}
\item \textbf{Byggföretag} – karaktäriseras av förekomst av osund konkurrens, arbetskraftsrelaterad brottslighet och systematiska fakturabedrägerier.
\item \textbf{Städföretag} – ofta beroende av lågavlönad arbetskraft och kontroller som är svåra att upprätthålla, vilket ökar risken för skatteundandragande och svart arbetskraft.
\item \textbf{Restaurangverksamheter} – ofta kontantintensiva och benägna till skatteundandragande, oredovisade löner samt manipulering av kassaflöden.
\item \textbf{Transportverksamheter} – särskilt i underleverantörsledet noteras bristfälliga kontroller, svarta löner och utrymme för fiktiv fakturering.
\end{itemize}

Mot bakgrund av såväl myndighetsanalysen som vår operativa erfarenhet fastställer vi hotnivån till \textbf{4 – hög extern risk} för dessa kundtyper.



%%%%%%%%%%%%%%%%%%%%%%%%%%%%%%%%%%%%%%%%%%%%%%%%%%%%%%%%%%%%%%%%%%%%%%%%%%%%%%%%%%%%%%%%%%%%%%%%%%%%%%%%%%%%%%%%%%%%%%%%%%%%%%%%%%%%%%%%%%%%%%%%%%%%%%%%%%


% Spårbarhet
\newpage
\vspace{3em}
\section*{Motivering av sårbarhetsnivåer per tjänst}
\addcontentsline{toc}{section}{Motivering av sårbarhetsnivåer per tjänst}

Vi väljer i denna analys att klassificera samtliga våra tjänster med sårbarhetsnivån \textbf{3 – betydande}, i linje med slutsatserna i den myndighetsgemensamma rapporten från 2020/2021. Tjänsterna bygger i stor utsträckning på kundens egna uppgifter och är därmed exponerade för manipulation, även om våra interna rutiner syftar till att motverka detta. 

Vi tillämpar denna nivå som en försiktighetsåtgärd och är öppna för revidering vid förändrade förutsättningar eller vägledning från tillsynsmyndighet.


\subsection*{1. Löpande bokföring – Sårbarhet = 3 (betydande)}
Enligt den nationella riskbedömningen (kap. 7.15, s. 91) är redovisningsbyråer en särskilt utsatt sektor eftersom de i praktiken kan möjliggöra eller maskera brott genom manipulativ bokföring. Tjänsten bygger på kundens underlag och är ofta beroende av kundens ärlighet. Vi bedömer därför sårbarheten som \textbf{betydande}.

\subsection*{2. Momsdeklaration – Sårbarhet = 3 (betydande)}
Momsfusk utgör enligt rapporten ett vanligt och välkänt modus för ekonomisk brottslighet (s. 40–41). Även om deklarationen i sig är formbunden, är underlaget mottagligt för manipulation. Därför klassificeras sårbarheten som \textbf{betydande}.

\subsection*{3. Årsbokslut/Årsredovisning – Sårbarhet = 3 (betydande)}
Årsbokslut och årsredovisningar är till stor del sammanställningar av redan hanterade transaktioner, men enligt rapporten (s. 85) används de i hög grad för att vilseleda banker och myndigheter. Detta innebär en betydande sårbarhet eftersom manipulation ofta sker genom selektiv presentation av data.

\subsection*{4. Inkomstdeklaration – Sårbarhet = 3 (betydande)}
Trots att deklarationen skickas till Skatteverket sker själva datainsamlingen internt eller via konsult. Deklarationen är därför beroende av kundens rapportering och dokumentation, vilket öppnar för otillbörlig påverkan. Riskbedömningen stöder att sårbarheten för denna typ av tjänst bedöms som \textbf{betydande}.





%%%
% Exempel Hot x Sårbarhet

\vspace{3em}
\subsection*{Exempel på tillämpning av riskmatrisen}
\addcontentsline{toc}{subsection}{Exempel på tillämpning av riskmatrisen}

För att tydliggöra hur vår riskbedömning omsätts i praktiken ger vi här ett konkret exempel. Vi visar hur våra tjänster relaterar till olika kundtyper utifrån två axlar: hot och sårbarhet.

Varje kombination ger upphov till en risknivå genom att multiplicera de två värdena. Den resulterande risken (Hot $\times$ Sårbarhet) ligger till grund för vår uppföljning, dokumentation och val av rutiner.

\begin{center}
\renewcommand{\arraystretch}{1.5}
\begin{tabular}{|p{5cm}|c|c|}
\hline
\textbf{Tjänst / Kundtyp} 
& \shortstack{\textbf{Kommunalt bolag} \\ \textbf{(Hot 1)}} 
& \shortstack{\textbf{Restaurang} \\ \textbf{(Hot 4)}} \\
\hline
\textbf{Löpande bokföring (Sårbarhet 3)} & $1 \times 3 = 3$ & $4 \times 3 = 12$ \\
\hline
\textbf{Momsdeklaration (3)} & $1 \times 3 = 3$ & $4 \times 3 = 12$ \\
\hline
\textbf{Årsbokslut (3)} & $1 \times 3 = 3$ & $4 \times 3 = 12$ \\
\hline
\textbf{Inkomstdeklaration (3)} & $1 \times 3 = 3$ & $4 \times 3 = 12$ \\
\hline
\end{tabular}
\end{center}



\vspace{1em}

\subsection*{Kommentar till metod och tillämpning}

Denna tabell visar hur risknivån för samma tjänst varierar beroende på kundens verksamhetstyp. Exempelvis ger en restaurangkund som vi bistår med inkomstdeklaration en riskpoäng på 12 (Hot = 3, Sårbarhet = 4), vilket klassificeras som \textit{hög risk}. I ett sådant fall krävs ytterligare åtgärder enligt vår interna rutin, såsom förstärkt kundkännedom och manuell granskning av verifikat från kontantkassa.

Däremot ger samma tjänst till ett kommunalt fastighetsbolag en riskpoäng på endast 4, vilket klassificeras som \textit{låg till medelrisk} och motiverar standardrutiner.








\end{document}
